\chapter{About this book}

This mini-book is a print adaptation of the Physical
Security framework from the \textbf{Security Frameworks},
an initiative by the \textbf{Security Alliance}. The web
pages were written for a site. This edition is rewritten
for a pocket book you can hand to someone.

The source this edition follows is Coercion \& Duress:
the reviewed core of the current framework. Facility
security, counter-surveillance, and hardware supply-chain
integrity still exist online as stubs. They are not in
this book. Empty chapters make a worse companion than a
short one.

This print cut changes order, cuts website navigation,
and turns the same controls into reading you can do on a
train. It does not invent new controls to fill gaps.

\section{Help us improve}

The live framework moves. If a control here is wrong,
incomplete, or now dangerous in your jurisdiction, fix it
upstream so the site and the next printing match.

{\ttfamily github.com/\\security-alliance/frameworks}

\keepblock{%
If you value this kind of resource, you can support the
work at our website.

{\centering\plainurl{securityalliance.org/donate}\par}
}

\keepblock{%
\section{Human safety comes first}

No asset, key, policy, or evidence goal is worth more
harm. If someone is being coerced, the job is to keep
them alive and to stop the situation getting worse.
Technical containment starts only when it does not put
the victim, family, staff, responders, or bystanders in
more danger.
}

\chapter{Introduction}

Cryptography protects keys from math. It does not protect
the person who can use those keys. Wrench attacks skip
the protocol and go after people: threats, kidnapping,
extortion, forced signing, raids.

Self-custody and on-chain treasuries put large balances
on a small number of humans. Public posts, conference
badges, and leaked addresses make those humans easier to
find. Most wallet guidance still talks as if the only
failure is a stolen seed on a laptop.

This book is for founders, traders, operators, signers,
and the people around them. It is also for the org that
will have to respond if one of those people is grabbed.

Read it as a loop, not as a policy you recite once:

\begin{enumerate}
  \item See who can reach you, and what one person can
  move.
  \item Design keys and roles so coercing one person is
  not enough.
  \item Shrink what follows you around: devices, family,
  routines, public identity.
  \item If it happens, protect life first. Move funds
  only when that does not raise the violence.
  \item Afterward, look after people, throw away exposed
  material, and close the hole that was used.
\end{enumerate}

\section{Disclaimer}

This is not legal advice, and it is not a close-protection
course. It will not match every country, family, or
treasury. Some of the dollar figures and wallet splits
in later chapters are examples from the source, not
defaults. Copying them blindly is how you get a setup
an attacker can predict.

A hidden duress path, a fake-empty wallet, or a
``surrender'' balance can get someone hurt if it is
obviously fake, or if it contradicts what the attacker
already saw on-chain. Design that kind of control with
people who do this work. Do not invent it in a parking
lot.

The contents are free on the Security Frameworks site.
Suggest changes there.

\section{How this book is organized}

\begin{enumerate}
  \item \textbf{The threat.} What a wrench attack is,
  who does it, what they want, where they hit.
  \item \textbf{How they pick you.} Visibility, routines,
  concentrated signing power, early signals.
  \item \textbf{Design the keys.} Minimization, split
  roles, delay, a believable low-value path.
  \item \textbf{Protect the people.} What follows a
  signer around, including family and staff.
  \item \textbf{If it happens.} Life first, then limits,
  then containment when it is safe.
  \item \textbf{After.} Support, rebuild, close the gap.
  \item \textbf{Do this next.} A phased list so this does
  not stay theoretical.
\end{enumerate}

The last chapter is four short cards for the moment you
do not want to reread a section.
